\documentclass[11pt,a4paper]{article}
\usepackage[utf8]{inputenc}
\usepackage[T1]{fontenc}
\usepackage[margin=2.4cm]{geometry}
\usepackage{lmodern}
\usepackage{microtype}

\pagestyle{empty}

\begin{document}

\begin{flushleft}
    {\small\textsc{M.Sc. Project Proposal --- Social Robotics}} \\[6pt]
    {\LARGE\bfseries An Ontology-Grounded Knowledge Graph Architecture for Explainable Robot Emotional Modeling and Expressive 3D Interaction} \\[8pt]
    {\small Based on \textit{Case Study 2: Ontology-Based Robot Emotional Modeling} (Miccoli et al., IAxIA 2026)} \\[4pt]
    {\small Repository: \texttt{https://github.com/Step0491/Social-Robotics-Emotional-Robot}}
\end{flushleft}

\vspace{4pt}
\noindent\rule{\textwidth}{0.5pt}
\vspace{8pt}

\section*{Abstract}
Social robots operating in human-centered environments require transparent and consistent mechanisms to model their internal affective states and adapt their behavioral responses across diverse interaction contexts. While Large Language Models (LLMs) enable fluent dialogue generation, relying on them for affective decision-making introduces opacity, lack of contextual guarantees, and limited behavioral consistency. Building upon an ontology-grounded Knowledge Graph framework for explainable emotional support, this project addresses \textbf{Ontology-Based Robot Emotional Modeling} by shifting the modeling perspective from the human user to the social robot's simulated internal emotional state. We design a formal domain ontology and a typed Knowledge Graph representing perceived social events, interaction contexts, internal emotional states, candidate behavioral responses, and contextual feasibility constraints. Decision-making is performed deterministically through graph-based reasoning, which updates the robot's emotional state, filters contraindicated behaviors, and produces an explicit \textit{why/why-not} reasoning trace. A locally deployed LLM (via Ollama) is strictly grounded in this reasoning trace to generate context-aware verbal responses without participating in the decision logic. Furthermore, the robot's structured emotional state and intensity are mapped onto an interactive \textbf{3D expressive robotic face}, enabling real-time visual feedback synchronized with the underlying symbolic reasoning. Finally, the system is evaluated across synthetic safe and constrained interaction scenarios in terms of behavioral coherence, constraint enforcement, and reasoning latency.

\vspace{6pt}

\section*{System Architecture \& Implementation Plan}
\begin{itemize}
    \item \textbf{Ontology \& Knowledge Graph (TBox/ABox):} Formal representation of robot emotions, perceived events, interaction contexts, behavioral responses, and feasibility constraints (\texttt{contraindicated\_in}), with episodic state tracking and OWL/RDF exportability.
    \item \textbf{Perception Extraction \& Validation Layer:} Structured extraction of events and contexts from natural language inputs via local LLM (Ollama), strictly validated against the domain ontology before updating the graph.
    \item \textbf{Explainable Graph Reasoning \& Grounded Generation:} Deterministic graph traversal for emotional state transitions, constraint filtering, \textit{why/why-not} reasoning traces, and grounded verbal response generation via Ollama.
    \item \textbf{3D Expressive Face \& Evaluation:} Real-time WebGL/Three.js 3D robotic face driven by affective parameters (valence, arousal, intensity) and quantitative validation on safe/constrained scenarios.
\end{itemize}

\end{document}
